\documentclass{article}

\usepackage{float}
\restylefloat{table}

\usepackage{booktabs}

\title{Team Productivity: POC\\\progname}

\author{\authname}

\date{}

\input{../Comments.text}
\input{../Common.text}

\begin{document}

\maketitle

This document summarizes the contributions of each team member up to the POC
Demo.  The time period of interest is the time between the beginning of the term
and the POC demo.

\wss{Please do not delete sections or modify the table formats.  Standardization
helps the instructors and TAs to review these documents. You may add sections to
the report.}

\wss{Remove the instructor guidelines/comments in the final version.}

\wss{Fill in the requested information, like project url, TA name, supervisor name, etc.}

\section{Demo Plans}

\wss{What is your feasibility risk?}

\wss{What will you be demonstrating}

\section{Team Meeting Attendance}
This section outlines the total number of group meetings each team member attended.

\begin{table}[H]
\centering
\begin{tabular}{ll}
\toprule
\textbf{Student} & \textbf{Meetings}\\
\midrule
Total & 10\\
Nour Fawaz & 10 \\
Javayria Mudassar & 10 \\
Yilin Jiang & 9 \\
Catherine Yang & 10 \\
\bottomrule
\end{tabular}
\end{table}

\begin{itemize}
    \item Yilin was unable to attend one team meeting due to illness.
\end{itemize}

\section{Supervisor/Stakeholder Meeting Attendance}

\noindent \textbf{Supervisor's Name: } Dr. Irene Ye Yuan

\begin{table}[H]
\centering
\begin{tabular}{ll}
\toprule
\textbf{Student} & \textbf{Meetings}\\
\midrule
Total & 5\\
Nour Fawaz & 5 \\
Javayria Mudassar & 4 \\
Yilin Jiang & 4 \\
Catherine Yang & 5 \\
\bottomrule
\end{tabular}
\end{table}
\begin{itemize}
    \item Javayria was unable to attend one supervisor meeting due to an injury (concussion).
    \item Yilin was unable to attend one supervisor meeting due to illness.
\end{itemize}

\section{Lecture Attendance}
This section outlines each group member’s lecture attendance from the beginning of classes up until the PoC demos.

\begin{table}[H]
\centering
\begin{tabular}{ll}
\toprule
\textbf{Student} & \textbf{Lectures}\\
\midrule
Total & 11\\
Nour Fawaz & 9 \\
Javayria Mudassar & 6 \\
Yilin Jiang & 6 \\
Catherine Yang & 7 \\
\bottomrule
\end{tabular}
\end{table}
\noindent
Although not every member attended every lecture, the team ensured that at least one member was present at each one so they could update the others afterward.


\section{TA Document Discussion Attendance}
This section records each group member’s attendance at all TA document discussion meetings, from the beginning of classes up until the PoC demos.

\noindent \textbf{TA's Name: } Tanya Djavaherpour

\begin{table}[H]
\centering
\begin{tabular}{ll}
\toprule
\textbf{Student} & \textbf{Meetings}\\
\midrule
Total & 4\\
Nour Fawaz & 4 \\
Javayria Mudassar & 3 \\
Yilin Jiang & 4 \\
Catherine Yang & 4 \\
\bottomrule
\end{tabular}
\end{table}

\begin{itemize}
    \item Javayria was unable to attend the supervisor meeting due to an injury (concussion).
\end{itemize}

\section{Commits, Lines Added and Lines Deleted, Pull Requests}

\wss{For each team member how many commits to the main branch have been made
over the time period of interest.  The total is the total number of commits for
the entire team since the beginning of the term.  The percentage is the
percentage of the total commits made by each team member.}

\wss{There are also columns for lines added and lines deleted to give a sense of
the magnitude of each team member's contributions.  Note that these numbers can
be somewhat misleading, as a team member who does refactoring may have high
numbers in both columns, while a team member who implements a large new feature
may have a high number of lines added, but few lines deleted.  As for all commit
information, these numbers will be interpreted with care.}

\wss{The pull requests column counts the number of pull requests created by each
team member.  This is a measure of how many contributions were made in branches,
which may not be reflected in the commit counts for the main branch.}

\begin{table}[H]
\centering
\begin{tabular}{lllllll}
\toprule
\textbf{Student} & \textbf{Commits} & \textbf{Percent} & \textbf{Lines Added} &
\textbf{Lines Deleted} & \textbf{Pull Requests}\\
\midrule
Total & Num & 100\% & total added & total deleted & total PRs\\
Name 1 & Num & \% & num & num & num\\
Name 2 & Num & \% & num & num & num\\
Name 3 & Num & \% & num & num & num\\
Name 4 & Num & \% & num & num & num\\
Name 5 & Num & \% & num & num & num\\
\bottomrule
\end{tabular}
\end{table}

\wss{If needed, an explanation for the counts can be provided here.  For
instance, if a team member has more commits to unmerged branches, these numbers
can be provided here.  If multiple people contribute to a commit, git allows for
multi-author commits.}

\section{Issue Tracker}

\wss{For each team member how many issues have they authored (including open and
closed issues (O+C)) and how many have they been assigned (only counting closed
issues (C only)) over the time period of interest.}

\begin{table}[H]
\centering
\begin{tabular}{lll}
\toprule
\textbf{Student} & \textbf{Authored (O+C)} & \textbf{Assigned (C only)}\\
\midrule
Name 1 & Num & Num \\
Name 2 & Num & Num \\
Name 3 & Num & Num \\
Name 4 & Num & Num \\
Name 5 & Num & Num \\
\bottomrule
\end{tabular}
\end{table}

\wss{If needed, an explanation for the counts can be provided here.}

\section{CICD via GitHub Actions}

\wss{Say how CICD technology is used and will be used in your project.}

\wss{Provide links to your CICD yaml files}

\section{Custom Documents (CDs)}

\wss{What is the plan (as documented in TeamComposition.csv) for the team's
CDs?  Should the CDs be modified now that the team knows more about the
project?}

\section{Team Charter Trigger Items}

\wss{Provide a summary of the quantified triggers identified in the team's
charter.}

\wss{Provide a list of any violations of the triggers.  If the team wishes, the
violations can be summarized on aggregate, instead of naming specific team
members.}

\wss{Provide a plan to address the violations.  This could include revising the
triggers, if they are found to be too weak, strong or ambiguous.}

\section{Additional Productivity Metrics}

\wss{If your team has additional metrics of productivity, please feel free to
add them to this report.  If not, please explicitly state that there are no additional
metrics.}

\wss{Additional metrics can include things like code reviews done, pull requests
created, count of joining meetings late, count of number of times contributions
had to be corrected, number of internal deadlines missed, test cases written, etc.}

\wss{We are looking for data on these metrics, not just a list of additional
metrics the team is planning on tracking.  However, if all you have is a plan,
please share it here.}

\end{document}